% AI Destekli 3 Katmanlı Test Otomasyon Framework'ü — Bitirme / Vize Raporu
% Derleme: pdflatex rapor.tex (iki kez — içindekiler için)
\documentclass[12pt,a4paper,twoside]{report}

\usepackage[utf8]{inputenc}
\usepackage[T1]{fontenc}
\usepackage[turkish]{babel}
\usepackage{geometry}
\usepackage{graphicx}
\usepackage{hyperref}
\usepackage{booktabs}
\usepackage{listings}
\usepackage{xcolor}
\usepackage{float}
\usepackage{placeins}
\usepackage{caption}
\usepackage{subcaption}
\usepackage{amsmath}
\usepackage{setspace}
\usepackage{titlesec}

\geometry{margin=2.5cm}
\onehalfspacing
\setlength{\textfloatsep}{8pt plus 2pt minus 4pt}
\setlength{\intextsep}{8pt plus 2pt minus 4pt}
\setlength{\floatpagefraction}{0.85}
\renewcommand{\topfraction}{0.9}
\renewcommand{\bottomfraction}{0.85}
\hypersetup{colorlinks=true, linkcolor=blue, urlcolor=blue, citecolor=blue}

\lstset{
  basicstyle=\ttfamily\small,
  breaklines=true,
  frame=single,
  backgroundcolor=\color{gray!8},
  keywordstyle=\color{blue},
  commentstyle=\color{green!50!black}
}

% --- Kapak bilgileri (üniversite şablonuna göre düzenle) ---
\newcommand{\raporBaslik}{AI Destekli 3 Katmanlı Web Test Otomasyon Framework'ü}
\newcommand{\raporAltBaslik}{Self-Healing, Yapay Zeka Destekli Hata Analizi ve QA Command Center Dashboard}
\newcommand{\ogrenciAd}{Ad Soyad}
\newcommand{\ogrenciNo}{Öğrenci No}
\newcommand{\danisman}{Danışman Adı Soyadı}
\newcommand{\universite}{Üniversite Adı}
\newcommand{\bolum}{Bölüm Adı}
\newcommand{\tarih}{Haziran 2026}

\begin{document}

% ===================== KAPAK =====================
\begin{titlepage}
  \centering
  \vspace*{1.5cm}
  {\Large \universite \par}
  {\large \bolum \par}
  \vspace{2cm}
  {\LARGE\bfseries \raporBaslik \par}
  \vspace{0.8cm}
  {\large \raporAltBaslik \par}
  \vspace{2.5cm}
  {\large Bitirme Projesi / Vize Raporu \par}
  \vspace{2cm}
  \begin{tabular}{rl}
    Öğrenci  & : \ogrenciAd \\
    Numara   & : \ogrenciNo \\
    Danışman & : \danisman \\
    Tarih    & : \tarih \\
  \end{tabular}
  \vfill
\end{titlepage}

% ÖZET ve ABSTRACT — güncel proje durumuna göre (Haziran 2026)

\phantomsection
\addcontentsline{toc}{section}{ÖZET}
\begin{center}
\textbf{\large ÖZET}\\[0.4em]
\textbf{E-TİCARET UYGULAMALARI İÇİN SELF-HEALING VE YAPAY ZEKA DESTEKLİ\\
ÜÇ KATMANLI TEST OTOMASYON FRAMEWORK'Ü}
\end{center}

Bu çalışma kapsamında, e-ticaret web uygulamalarının kalite güvencesini sağlamak amacıyla sürdürülebilir, genişletilebilir ve yapay zeka destekli bir test otomasyon framework'ü geliştirilmiştir. Projenin temel amacı, manuel test süreçlerinin getirdiği zaman ve maliyet yükünü azaltmak; UI değişikliklerine karşı dayanıklı, veri odaklı ve standartlara dayalı bir otomasyon altyapısı kurmaktır.

Framework, Java~17, Selenium WebDriver~4, TestNG, Maven, RestAssured ve JDBC teknolojileri üzerine inşa edilmiştir. Mimari olarak Page Object Model (POM) deseni benimsenmiş; test verileri JSON dosyalarında tutularak TestNG DataProvider mekanizmasıyla teste aktarılmıştır. Bu sayede test kodu ile test verisi birbirinden ayrıştırılmış, bakım maliyeti minimize edilmiştir. Proje yalnızca arayüz katmanıyla sınırlı kalmayıp \textbf{UI, API, veritabanı (DB) ve uçtan uca (E2E)} olmak üzere dört test katmanını tek çatı altında birleştirmektedir.

Projenin teknik ayırt ediciliği üç bileşende yoğunlaşmaktadır. Birincisi, UI değişikliklerine karşı dayanıklılık sağlayan \textbf{self-healing} mekanizmasıdır. Her element için öncelikli locator listesi tanımlanmış; birincil locator başarısız olduğunda sistem otomatik olarak yedek locator'ları denemekte ve kullanım istatistiklerini raporlamaktadır. İkincisi, test hatalarının kök neden analizini otomatikleştiren \textbf{LLM tabanlı failure analysis} bileşenidir. Test başarısız olduğunda hata mesajı, modül metadata'sı, self-healing geçmişi ve ekran görüntüsü otomatik toplanarak Groq LLM servisine iletilmekte; model Türkçe kök neden, önerilen aksiyon ve risk değerlendirmesi üretmektedir. Opsiyonel Flask tabanlı makine öğrenmesi servisi ise flaky/stabil test tahminini desteklemektedir. Üçüncüsü, koşum sonuçlarını görselleştiren \textbf{QA Command Center} web dashboard'udur; test geçmişi, modül trendi, OWASP eşlemesi, AI analiz notları ve Jira formatında hata raporlama akışını sunmaktadır.

Test senaryoları ISTQB test tasarım teknikleri (eşdeğerlik bölümlemesi, sınır değer analizi, negatif test, durum geçiş testi, güvenlik testi) ve OWASP güvenlik standartları baz alınarak tasarlanmıştır. Modül karmaşıklığı ve iş kritikliğine göre risk bazlı senaryo hedefleri belirlenmiştir. Framework kapsamında Login, Register, Forgot Password, My Account, My Account Settings, Cart, Checkout, Search, WishList ve Newsletter UI modülleri; RestAssured ile API katmanı; H2 ve Microsoft SQL Server ile DB katmanı; UI--API--DB entegrasyonunu kapsayan E2E sipariş akışı test edilmiştir. Toplam \textbf{273 data-driven test senaryosu} uygulanmıştır. Test sonuçları, self-healing telemetrisi, teknik hata logları, LLM analizleri ve ekran görüntüleri \texttt{target/test-history} dizininde toplanarak dashboard üzerinden görselleştirilmektedir.

\vspace{2cm}
\textbf{Anahtar Kelimeler:} Test Otomasyonu, Selenium, Self-Healing, Büyük Dil Modeli (LLM), ISTQB, OWASP, Page Object Model, Data-Driven Test, Üç Katmanlı Test, Web Dashboard, Java

\newpage

\phantomsection
\addcontentsline{toc}{section}{ABSTRACT}
\begin{center}
\textbf{\large ABSTRACT}\\[0.4em]
\textbf{A SELF-HEALING AND AI-POWERED THREE-LAYER\\
TEST AUTOMATION FRAMEWORK FOR E-COMMERCE APPLICATIONS}
\end{center}

In this study, a sustainable, extensible and AI-powered test automation framework has been developed to ensure the quality assurance of e-commerce web applications. The primary objective of the project is to reduce the time and cost burden of manual testing processes and to establish an automation infrastructure that is resilient to UI changes, data-driven and standards-based.

The framework is built on Java~17, Selenium WebDriver~4, TestNG, Maven, RestAssured and JDBC technologies. The Page Object Model (POM) pattern has been adopted as the architectural approach; test data is stored in JSON files and transferred to tests via the TestNG DataProvider mechanism. This approach separates test code from test data, minimizing maintenance costs. The project extends beyond the UI layer by integrating four test layers under a single framework: \textbf{UI, API, database (DB) and end-to-end (E2E)}.

The technical distinction of the project focuses on three components. The first is the \textbf{self-healing} mechanism that provides resilience against UI changes. A prioritized locator list is defined for each element; when the primary locator fails, the system automatically tries fallback locators and records usage telemetry. The second is the \textbf{LLM-based failure analysis} component that automates root cause analysis of test failures. When a test fails, the error message, module metadata, self-healing history and screenshot are automatically collected and sent to the Groq LLM service, which generates root cause analysis, recommended actions and risk assessment in Turkish. An optional Flask-based machine learning service supports flaky/stable test prediction. The third is the \textbf{QA Command Center} web dashboard, which provides test history, module trends, OWASP mapping, AI analysis notes and a Jira-style bug reporting workflow.

Test scenarios have been designed based on ISTQB test design techniques (equivalence partitioning, boundary value analysis, negative testing, state transition testing, security testing) and OWASP security standards. Risk-based scenario targets have been defined according to module complexity and business criticality. Within the framework, Login, Register, Forgot Password, My Account, My Account Settings, Cart, Checkout, Search, WishList and Newsletter UI modules; the API layer with RestAssured; the DB layer with H2 and Microsoft SQL Server; and the E2E order flow integrating UI--API--DB have been tested. A total of \textbf{273 data-driven test scenarios} have been implemented. Test results, self-healing telemetry, technical failure logs, LLM analyses and screenshots are collected under \texttt{target/test-history} and visualized through the dashboard.

\vspace{2cm}
\textbf{Keywords:} Test Automation, Selenium, Self-Healing, Large Language Model (LLM), ISTQB, OWASP, Page Object Model, Data-Driven Testing, Three-Layer Testing, Web Dashboard, Java

\pagebreak{}


\pagenumbering{roman}
\tableofcontents
\clearpage
\listoffigures
\clearpage
\listoftables
\clearpage

% Manuel içindekiler (tahmini sayfa numaraları) — otomatik TOC ile birlikte veya yerine kullanılabilir
% Manuel içindekiler — tahmini sayfa numaraları (içerik uzunluğuna göre güncelle)
% Otomatik \tableofcontents ile birlikte kullanılabilir veya yorum satırına alınabilir.

\chapter*{İçindekiler (Tahmini Sayfa Numaraları)}
\addcontentsline{toc}{chapter}{İçindekiler (Tahmini Sayfa Numaraları)}

\newcommand{\tocentry}[2]{%
  \noindent\parbox[t]{0.82\textwidth}{#1}\hfill\parbox[t]{0.12\textwidth}{\raggedleft #2}\par\vspace{0.35em}}
\newcommand{\tocsub}[2]{%
  \noindent\hspace{1.5em}\parbox[t]{0.78\textwidth}{#1}\hfill\parbox[t]{0.12\textwidth}{\raggedleft #2}\par\vspace{0.25em}}
\newcommand{\tocsubsub}[2]{%
  \noindent\hspace{3em}\parbox[t]{0.74\textwidth}{#1}\hfill\parbox[t]{0.12\textwidth}{\raggedleft #2}\par\vspace{0.2em}}

\tocentry{\textbf{Önsöz}}{iii}
\tocentry{\textbf{Kısaltmalar ve Semboller}}{iv}
\par\vspace{0.5em}

\tocentry{\textbf{1. Giriş ve Motivasyon}}{1}
\tocsub{1.1. Problemin Tanımı}{1}
\tocsub{1.2. Projenin Amacı ve Kapsamı}{2}
\tocsub{1.3. Hedeflenen Katmanlar (UI, API, DB, E2E)}{3}
\tocsub{1.4. Raporun Organizasyonu}{4}

\tocentry{\textbf{2. Kullanılan Teknolojiler}}{5}
\tocsub{2.1. Java 17 ve Maven}{5}
\tocsub{2.2. Selenium WebDriver 4}{6}
\tocsub{2.3. TestNG ve ExtentReports}{7}
\tocsub{2.4. RestAssured (API Katmanı)}{8}
\tocsub{2.5. JDBC — H2 ve Microsoft SQL Server}{9}
\tocsub{2.6. Jackson (JSON Data-Driven)}{10}
\tocsub{2.7. Groq LLM ve Flask ML Servisi}{11}

\tocentry{\textbf{3. Sistem Mimarisi}}{13}
\tocsub{3.1. Genel Mimari Görünüm}{13}
\tocsub{3.2. Page Object Model (POM)}{15}
\tocsub{3.3. ConfigManager ve Ortam Yönetimi}{16}
\tocsub{3.4. TestNG Listener ve Raporlama Hattı}{17}
\tocsub{3.5. DashboardResultsListener ve test-history}{18}
\tocsub{3.6. Üç Katman + E2E Entegrasyon Akışı}{20}

\tocentry{\textbf{4. Uygulanan Test Yaklaşımı}}{22}
\tocsub{4.1. Data-Driven Test Tasarımı (JSON Senaryolar)}{22}
\tocsub{4.2. Risk Bazlı Senaryo Hedefleri}{24}
\tocsub{4.3. Test Teknikleri (EP, BVA, Negatif, Güvenlik)}{25}
\tocsub{4.4. OWASP ve ISTQB Standart Eşlemesi}{26}
\tocsub{4.5. UI Modül Kapsamı (~202 senaryo)}{27}
\tocsub{4.6. API, DB ve E2E Senaryo Kapsamı}{29}

\tocentry{\textbf{5. Self-Healing Mekanizması}}{31}
\tocsub{5.1. Motivasyon ve Problem}{31}
\tocsub{5.2. SelfHealingElementFinder Tasarımı}{32}
\tocsub{5.3. Yedek Locator Zinciri}{33}
\tocsub{5.4. Telemetri ve Dashboard Entegrasyonu}{34}
\tocsub{5.5. Örnek Kullanım Senaryosu}{35}

\tocentry{\textbf{6. Yapay Zeka Destekli Hata Analizi}}{37}
\tocsub{6.1. AiFailureAnalyzer Mimarisi}{37}
\tocsub{6.2. Groq LLM ile Kök Neden Analizi}{38}
\tocsub{6.3. Flask ML Modeli (Flaky/Stabil Tahmin)}{39}
\tocsub{6.4. FailureLogWriter — Teknik Hata Logları}{40}
\tocsub{6.5. Türkçe Analiz Çıktı Formatı}{41}

\tocentry{\textbf{7. QA Command Center Dashboard}}{43}
\tocsub{7.1. Dashboard Mimarisi (Tek HTML + Chart.js)}{43}
\tocsub{7.2. Genel Bakış ve Metrik Kartları}{44}
\tocsub{7.3. Test Geçmişi ve Modül Trendi}{45}
\tocsub{7.4. Hata Analizi ve AI Notları}{46}
\tocsub{7.5. Modül Kapsamı ve OWASP Haritası}{47}
\tocsub{7.6. Raporlanan Hatalar (Master-Detail)}{48}
\tocsub{7.7. results.json Veri Şeması}{49}

\tocentry{\textbf{8. Jira Hata Raporlama Akışı}}{51}
\tocsub{8.1. JiraBugReportBuilder}{51}
\tocsub{8.2. Türkçe Rapor Şablonu}{52}
\tocsub{8.3. Ekran Görüntüsü ve Log Ekleri}{53}
\tocsub{8.4. localStorage ve jira-reports.json}{54}

\tocentry{\textbf{9. Mevcut Durum ve Sonuçlar}}{55}
\tocsub{9.1. Toplam Test Kapsamı (~273 senaryo)}{55}
\tocsub{9.2. Modül Bazlı Sonuç Özeti}{56}
\tocsub{9.3. Bilinçli Güvenlik Test Fail'leri}{57}
\tocsub{9.4. Performans ve Paralel Koşum (TestNG)}{58}
\tocsub{9.5. Karşılaşılan Zorluklar ve Çözümler}{59}
\tocsub{9.6. Proje Değerlendirmesi}{60}

\tocentry{\textbf{10. Gelecek Planlar}}{62}
\tocsub{10.1. CI/CD (GitHub Actions)}{62}
\tocsub{10.2. Jira REST API Entegrasyonu}{63}
\tocsub{10.3. Dashboard Dağıtım Otomasyonu}{63}
\tocsub{10.4. Test Verisi ve Flaky Test Yönetimi}{64}

\par\vspace{0.8em}
\tocentry{\textbf{Ek A: Test Modül Listesi}}{65}
\tocentry{\textbf{Ek B: Örnek Konfigürasyon}}{67}
\tocentry{\textbf{Kaynakça}}{69}

\clearpage

\pagenumbering{arabic}

% ===================== BÖLÜMLER =====================
\chapter{Giriş ve Motivasyon}

\section{Problemin Tanımı}
Modern web uygulamalarında kullanıcı arayüzü, REST API ve veritabanı katmanları birbirine sıkı bağımlılık içindedir. Manuel regresyon testleri tekrarlanabilir, ölçeklenebilir ve hızlı geri bildirim sağlayamaz. Ayrıca DOM değişiklikleri, kırılgan locator'lar ve flaky testler otomasyon projelerinin sürdürülebilirliğini zorlaştırır.

Bu proje, LambdaTest E-Commerce Playground üzerinde çalışan ve \textbf{UI · API · DB · E2E} katmanlarını tek çatı altında toplayan bir test otomasyon framework'ü geliştirmeyi hedefler.

\section{Projenin Amacı ve Kapsamı}
Projenin temel amacı:
\begin{itemize}
  \item Data-driven JSON senaryoları ile modüler UI testleri yazmak,
  \item API ve veritabanı katmanlarında bağımsız doğrulama yapmak,
  \item Self-healing locator ile DOM değişikliklerine dayanıklılık sağlamak,
  \item Groq LLM ve opsiyonel ML modeli ile hata kök neden analizi üretmek,
  \item QA Command Center dashboard ile test geçmişini görselleştirmek.
\end{itemize}

\section{Hedeflenen Katmanlar}
\begin{table}[H]
  \centering
  \caption{Proje katmanları ve teknolojiler}
  \begin{tabular}{@{}lll@{}}
    \toprule
    \textbf{Katman} & \textbf{Araç} & \textbf{Hedef} \\
    \midrule
    UI & Selenium 4 + TestNG & 10 modül, ~202 senaryo \\
    API & RestAssured & reqres.in, 32 senaryo \\
    DB & JDBC (H2 + SQL Server) & 35 senaryo \\
    E2E & UI + API + DB & Sipariş akışı, 4 senaryo \\
    \bottomrule
  \end{tabular}
\end{table}

\section{Raporun Organizasyonu}
Bu rapor; teknoloji seçimleri, sistem mimarisi, test yaklaşımı, self-healing, AI analiz, dashboard, Jira raporlama, sonuçlar ve gelecek planları bölümlerinden oluşmaktadır.

% TODO: Üniversite önsözü, teşekkür ve kısaltmalar bölümünü şablona göre ekle.

\chapter{Kullanılan Teknolojiler}

\section{Java 17 ve Maven}
Proje JDK 17 ile geliştirilmiştir. Maven bağımlılık yönetimi, derleme ve TestNG suite koşumları için kullanılır. Ana artefakt: \texttt{ai-3layer-test-framework}.

\section{Selenium WebDriver 4}
UI otomasyonunun temelini Selenium 4 oluşturur. WebDriverManager sürücü indirmesini otomatikleştirir. Page Object Model (POM) ile sayfa sınıfları ayrıştırılmıştır.

\section{TestNG ve ExtentReports}
TestNG paralel sınıf koşumu (\texttt{parallel="classes"}, \texttt{thread-count="4"}) ve data provider desteği sağlar. ExtentReports HTML raporu \texttt{target/reports/extent-report.html} konumuna yazılır.

\section{RestAssured (API Katmanı)}
API testleri RestAssured kütüphanesi ile \texttt{reqres.in} üzerinde yürütülür. JSON path doğrulama, HTTP durum kodu ve yanıt süresi kontrolleri yapılır.

\section{JDBC — H2 ve Microsoft SQL Server}
\begin{itemize}
  \item \textbf{H2:} Bellek içi veritabanı; hızlı DB senaryoları ve CI uyumluluğu.
  \item \textbf{SQL Server:} Gerçek ortam simülasyonu; sipariş doğrulama ve E2E akışı.
\end{itemize}

\section{Jackson (JSON Data-Driven)}
Test senaryoları \texttt{src/test/resources/testdata/*-scenarios.json} dosyalarında tutulur. \texttt{JsonDataLoader} ile TestNG data provider'a aktarılır.

\section{Groq LLM ve Flask ML Servisi}
\begin{itemize}
  \item \textbf{Groq API:} Fail anında Türkçe kök neden analizi (\texttt{AiFailureAnalyzer}).
  \item \textbf{Flask (\texttt{ai-analysis/}):} Opsiyonel ML modeli; flaky/stabil tahmin (\texttt{qa\_model.pkl}).
\end{itemize}

% TODO: Teknoloji seçim gerekçelerini tablo halinde genişlet.

\chapter{Sistem Mimarisi}

\section{Genel Mimari Görünüm}
Framework dört ana bileşenden oluşur: test katmanları (UI/API/DB/E2E), yardımcı servisler (self-healing, AI), dinleyiciler (listener) ve dashboard veri hattı.

\begin{figure}[H]
  \centering
  \fbox{\parbox{0.88\textwidth}{\centering\vspace{2cm}
    \textbf{[Şekil: Mimari diyagram]}\\[0.5em]
    TestNG $\rightarrow$ BaseTest $\rightarrow$ Page/API/DB $\rightarrow$ Listener $\rightarrow$ test-history\\
    \small Dashboard HTML $\leftarrow$ results.json + screenshots + logs
  \vspace{2cm}}}
  \caption{Genel sistem mimarisi (placeholder — draw.io / TikZ ile güncelle)}
  \label{fig:mimari}
\end{figure}

\section{Page Object Model (POM)}
Her UI modülü için ayrı page sınıfı (\texttt{LoginPage}, \texttt{CheckoutPage} vb.) tanımlanmıştır. Test sınıfları yalnızca iş akışını orchestrate eder; DOM etkileşimleri page katmanındadır.

\section{ConfigManager ve Ortam Yönetimi}
\texttt{config.properties} dosyası URL, kimlik bilgisi ve AI anahtarlarını tutar. \texttt{ConfigManager} ortam değişkeni override destekler (\texttt{AI\_API\_KEY}, \texttt{BASE\_URL} vb.).

\section{TestNG Listener ve Raporlama Hattı}
\texttt{TestListener} ExtentReports ve ekran görüntüsü yakalamayı yönetir. \texttt{DashboardResultsListener} her test sonucunu \texttt{results.json} formatına yazar.

\section{DashboardResultsListener ve test-history}
Koşum sonrası şu dosyalar üretilir:
\begin{itemize}
  \item \texttt{target/test-history/results.json}
  \item \texttt{target/test-history/screenshots/*.png}
  \item \texttt{target/test-history/logs/*.txt}
\end{itemize}

\section{Üç Katman + E2E Entegrasyon Akışı}
E2E senaryolarda UI üzerinden checkout tamamlanır, API ile senkronizasyon doğrulanır, SQL Server'da sipariş kaydı sorgulanır (\texttt{E2EOrderFlowTest}).

% TODO: Sequence diagram ekle (checkout E2E akışı).

\chapter{Uygulanan Test Yaklaşımı}

\section{Data-Driven Test Tasarımı (JSON Senaryolar)}
Her modül için JSON dosyasında senaryo tanımları bulunur. Örnek alanlar: \texttt{caseId}, \texttt{description}, \texttt{technique}, \texttt{standard}, \texttt{testType}, \texttt{expectedOutcome}.

\begin{lstlisting}[caption={Örnek senaryo alanları (login-scenarios.json)}]
{
  "caseId": "LGN-001",
  "description": "Gecersiz sifre ile giris",
  "technique": "negative-test",
  "standard": "OWASP",
  "testType": "regression",
  "testLevel": "system-test"
}
\end{lstlisting}

\section{Risk Bazlı Senaryo Hedefleri}
Modül karmaşıklığı ve iş kritikliğine göre hedef senaryo sayıları belirlenmiştir:

\begin{table}[H]
  \centering
  \caption{Risk bazlı modül hedefleri}
  \begin{tabular}{@{}lr@{}}
    \toprule
    \textbf{Modül} & \textbf{Hedef} \\
    \midrule
    Checkout & 28 \\
    Login / ForgotPassword & 24 \\
    Register / Cart / MyAccountSettings & 22 \\
    MyAccount & 20 \\
    Search & 18 \\
    WishList & 12 \\
    Newsletter & 10 \\
    Api & 32 \\
    Db (H2 + SQL Server) & 35 \\
    E2EOrderFlow & 4 \\
    \midrule
    \textbf{Toplam} & \textbf{$\sim$273} \\
    \bottomrule
  \end{tabular}
\end{table}

\section{Test Teknikleri}
Equivalence Partitioning (EP), Boundary Value Analysis (BVA), negatif/pozitif test, durum geçişi ve güvenlik odaklı senaryolar uygulanmıştır.

\section{OWASP ve ISTQB Standart Eşlemesi}
Dashboard Modül Kapsamı bölümünde OWASP test kategorileri modüllere eşlenmiştir. Senaryo metadata'sında \texttt{standard} alanı ile referans tutulur.

\section{UI Modül Kapsamı}
10 UI modülü: Login, Register, ForgotPassword, MyAccount, MyAccountSettings, Newsletter, WishList, Search, Cart, Checkout. Toplam $\sim$202 UI senaryosu.

\section{API, DB ve E2E Senaryo Kapsamı}
\begin{itemize}
  \item \textbf{API:} GET/POST/PUT/DELETE, auth, pagination, hata kodları (32 senaryo).
  \item \textbf{DB:} CRUD, constraint, transaction, SQL Server entegrasyon (35 senaryo).
  \item \textbf{E2E:} UI checkout $\rightarrow$ API $\rightarrow$ DB doğrulama (4 senaryo).
\end{itemize}

% TODO: Modül başına pass/fail tablosu ekle (son koşum sonuçlarından).

\chapter{Self-Healing Mekanizması}

\section{Motivasyon ve Problem}
Web uygulamalarında CSS seçicileri, ID'ler ve DOM yapısı sık değişir. Klasik otomasyon bu değişikliklerde kırılır. Self-healing, birincil locator başarısız olduğunda tanımlı yedek locator listesini sırayla dener.

\section{SelfHealingElementFinder Tasarımı}
\texttt{SelfHealingElementFinder} sınıfı, page object'lerde element bulma işlemini merkezileştirir. Başarılı yedek locator kullanıldığında telemetri kaydı oluşturulur.

\section{Yedek Locator Zinciri}
Her element için birden fazla \texttt{By} stratejisi tanımlanır (id, css, xpath, link text vb.). Zincir soldan sağa denenir; ilk görünür ve tıklanabilir element kabul edilir.

\section{Telemetri ve Dashboard Entegrasyonu}
Self-healing kullanımı \texttt{SelfHealingTelemetry} ile izlenir. Dashboard test kartlarında ``Otomatik Onarım Devrede'' rozeti gösterilir; \texttt{selfHealUsed}, \texttt{selfHealLocator} alanları \texttt{results.json}'a yazılır.

\section{Örnek Kullanım Senaryosu}
Checkout, Cart ve Search page sınıflarında self-healing aktiftir. Birincil sepet butonu locator'ı değiştiğinde yedek CSS seçicisi devreye girer ve test fail olmadan devam eder.

% TODO: Before/after locator tablosu ve gerçek self-heal log örneği ekle.

\chapter{Yapay Zeka Destekli Hata Analizi}

\section{AiFailureAnalyzer Mimarisi}
Test fail olduğunda \texttt{AiFailureAnalyzer.analyze()} çağrılır. İki katmanlı analiz uygulanır:
\begin{enumerate}
  \item Flask ML servisi: flaky/stabil tahmin (\texttt{/predict} endpoint).
  \item Groq LLM: Türkçe kök neden, önerilen aksiyon ve risk değerlendirmesi.
\end{enumerate}

\section{Groq LLM ile Kök Neden Analizi}
Groq OpenAI uyumlu API üzerinden \texttt{llama-3.1-8b-instant} modeli kullanılır. Analiz çıktısı şu başlıkları içerir:
\begin{itemize}
  \item Kök Neden
  \item URL / Assertion Detayları
  \item Önerilen Aksiyon
  \item Risk Değerlendirmesi
\end{itemize}

\section{Flask ML Modeli (Flaky/Stabil Tahmin)}
\texttt{ai-analysis/app.py} servisi geçmiş test metadata'sından eğitilmiş \texttt{qa\_model.pkl} modelini sunar. Modül, teknik, standart ve süre bilgileri feature olarak kullanılır.

\section{FailureLogWriter — Teknik Hata Logları}
Fail anında yalnızca teknik bilgi içeren log dosyası \texttt{target/test-history/logs/} altına yazılır. Stack trace, assertion mesajı ve test metadata'sı içerir.

\section{Türkçe Analiz Çıktı Formatı}
Dashboard Hata Analizi ve Raporlanan Hatalar bölümlerinde AI notu Türkçe gösterilir. \texttt{[AI-SKIPPED]} etiketi API anahtarı yoksa veya servis kapalıysa kullanılır.

% TODO: Gerçek FPW-EG-01 fail analizi örneğini listing olarak ekle.

\chapter{QA Command Center Dashboard}

\section{Dashboard Mimarisi (Tek HTML + Chart.js)}
Dashboard tek dosya olarak \texttt{dashboard.html} içinde geliştirilmiştir. Chart.js grafikleri, Lucide ikonları ve \texttt{results.json} fetch ile çalışır. Serve kökü: \texttt{target/test-history/}.

\section{Genel Bakış ve Metrik Kartları}
Ana sayfada toplam test, geçen, kalan ve başarı oranı kartları bulunur. Her kart drill-down popup açarak ilgili test listesini gösterir.

\section{Test Geçmişi ve Modül Trendi}
``Genel Koşumlar'' ve ``Modül Bazlı'' sekmeleri ile geçmiş koşumlar timeline'da listelenir. Modül trendi grouped bar chart ile modül bazlı pass/fail eğilimini gösterir.

\section{Hata Analizi ve AI Notları}
Fail olan testler katman filtresine göre (UI / API / DB) listelenir. Her kartta hata mesajı, ekran görüntüsü, AI analiz notu ve ``Hatayı Raporla'' butonu yer alır.

\section{Modül Kapsamı ve OWASP Haritası}
Modül kontrol listesi risk tier rozetleri ve dairesel ilerleme halkaları ile gösterilir. OWASP test kategorileri modüllere eşlenmiştir.

\section{Raporlanan Hatalar (Master-Detail)}
Jira tarzı master-detail düzen: sol listede raporlanan hatalar, sağda detay paneli. Ekler bölümünde ekran görüntüsü ve hata logu indirme desteklenir. Kayıtlar \texttt{localStorage} ile persist edilir.

\section{results.json Veri Şeması}
Her koşum \texttt{runs[]} dizisinde; testler \texttt{caseId}, \texttt{status}, \texttt{module}, \texttt{screenshotPath}, \texttt{errorLogPath}, \texttt{aiAnalysis} alanlarıyla saklanır.

% TODO: Dashboard ekran görüntüsü ekle (Genel Bakış, Hata Analizi).

\chapter{Jira Hata Raporlama Akışı}

\section{JiraBugReportBuilder}
Java tarafında \texttt{JiraBugReportBuilder} fail anında Türkçe Jira raporu üretir. Özet, test adımları, beklenen/gerçekleşen sonuç ve teknik metadata içerir.

\section{Türkçe Rapor Şablonu}
Dashboard Jira modal formu şu alanları sunar:
\begin{itemize}
  \item Proje anahtarı (varsayılan: \texttt{LTPLAY})
  \item Öncelik, önem, kanal, QA süreci
  \item Test adımları, beklenen sonuç, gerçekleşen sonuç
  \item Açıklama ve teknik hata detayı
\end{itemize}

\section{Ekran Görüntüsü ve Log Ekleri}
Fail anında screenshot \texttt{screenshots/} altına kopyalanır. \texttt{FailureLogWriter} teknik log üretir. Dashboard ``Ekleri İndir (.zip)'' ile her iki dosyayı birlikte indirmeyi destekler.

\section{localStorage ve jira-reports.json}
``Hatayı Raporla'' ile kayıt \texttt{qa-dashboard-reported-errors} localStorage anahtarına yazılır. Sunucu tarafı entegrasyon için \texttt{jira-reports.json} opsiyonel olarak kullanılabilir.

% TODO: Örnek Jira rapor metnini Ek C olarak ekle.

\chapter{Mevcut Durum ve Sonuçlar}

\section{Toplam Test Kapsamı}
Proje $\sim$273 data-driven senaryo içermektedir. UI regression suite (\texttt{ui-regression.xml}) 14 test sınıfını 6 TestNG grubunda koşturur.

\section{Modül Bazlı Sonuç Özeti}
% TODO: Son mvn clean test çıktısından güncel pass/fail sayılarını tabloya yaz.

\begin{table}[H]
  \centering
  \caption{Modül bazlı senaryo dağılımı (hedef)}
  \begin{tabular}{@{}lrr@{}}
    \toprule
    \textbf{Modül} & \textbf{Senaryo} & \textbf{Risk} \\
    \midrule
    Checkout & 28 & Yüksek \\
    Login & 24 & Yüksek \\
    ForgotPassword & 24 & Yüksek \\
    Register & 22 & Orta-Yüksek \\
    Cart & 22 & Orta-Yüksek \\
    MyAccountSettings & 22 & Orta-Yüksek \\
    MyAccount & 20 & Orta \\
    Search & 18 & Orta \\
    WishList & 12 & Düşük \\
    Newsletter & 10 & Düşük \\
    Api & 32 & — \\
    Db & 35 & — \\
    E2E & 4 & Yüksek \\
    \bottomrule
  \end{tabular}
\end{table}

\section{Bilinçli Güvenlik Test Fail'leri}
ForgotPassword modülünde OWASP odaklı senaryolar (\texttt{FPW-EG-01}, \texttt{FPW-ST-02}) bilinçli olarak fail tasarlanmıştır. Bu senaryolar gerçek güvenlik açıklarını test etmek içindir; framework hatası değildir.

\section{Performans ve Paralel Koşum}
TestNG \texttt{parallel="classes"} modu ile 4 thread kullanılır. ThreadLocal WebDriver ile sınıf düzeyinde izolasyon sağlanır.

\section{Karşılaşılan Zorluklar ve Çözümler}
\begin{itemize}
  \item \textbf{DOM değişiklikleri:} Self-healing locator zinciri.
  \item \textbf{Dashboard veri kaybı:} Modül koşumlarında \texttt{mvn clean} kullanılmaması.
  \item \textbf{Ek dosya yolları:} \texttt{resolveDashboardAsset()} ile path normalizasyonu.
  \item \textbf{Ürün katalog değişimi:} JSON senaryo güncellemesi (ör. WishList ürün adı).
\end{itemize}

\section{Proje Değerlendirmesi}
Proje; üç katmanlı otomasyon, self-healing, AI analiz ve görsel dashboard ile bitirme projesi kapsamını karşılamaktadır. Eksik kalan CI/CD ve Jira REST API entegrasyonu gelecek çalışma olarak planlanmıştır.

\chapter{Gelecek Planlar}

\section{CI/CD (GitHub Actions)}
GitHub Actions workflow ile her push'ta \texttt{mvn test} koşturulması, artefakt olarak ExtentReports ve \texttt{results.json} yüklenmesi planlanmaktadır.

\section{Jira REST API Entegrasyonu}
Dashboard formundan doğrudan Jira issue açma: \texttt{jira.enabled=true}, API token ve \texttt{JiraDashboardConfig} ile REST çağrısı.

\section{Dashboard Dağıtım Otomasyonu}
\texttt{dashboard.html} kopyalama ve \texttt{npx serve} adımlarını tek npm script veya Maven exec plugin ile otomatikleştirme.

\section{Test Verisi ve Flaky Test Yönetimi}
ML modelinin geçmiş koşum verisiyle periyodik yeniden eğitimi; flaky testlerin otomatik etiketlenmesi ve quarantine mekanizması.

% TODO: Zaman çizelgesi (Gantt) veya milestone tablosu ekle.


% ===================== EKLER =====================
\appendix
\chapter{Ek A: Test Modül Listesi}
\begin{table}[H]
  \centering
  \caption{UI test modülleri ve test sınıfları}
  \begin{tabular}{@{}lll@{}}
    \toprule
    \textbf{Modül} & \textbf{Test Sınıfı} & \textbf{JSON Dosyası} \\
    \midrule
    Login & LoginUiTest & login-scenarios.json \\
    Register & RegisterUiTest & register-scenarios.json \\
    ForgotPassword & ForgotPasswordUiTest & forgot-password-scenarios.json \\
    MyAccount & MyAccountUiTest & my-account-scenarios.json \\
    MyAccountSettings & MyAccountSettingsUiTest & my-account-settings-scenarios.json \\
    Newsletter & NewsletterUiTest & newsletter-scenarios.json \\
    WishList & WishListUiTest & wishlist-scenarios.json \\
    Search & SearchUiTest & search-scenarios.json \\
    Cart & CartUiTest & cart-scenarios.json \\
    Checkout & CheckoutUiTest & checkout-scenarios.json \\
    \bottomrule
  \end{tabular}
\end{table}

\begin{table}[H]
  \centering
  \caption{API, DB ve E2E test sınıfları}
  \begin{tabular}{@{}lll@{}}
    \toprule
    \textbf{Katman} & \textbf{Test Sınıfı} & \textbf{JSON Dosyası} \\
    \midrule
    API & ApiTest & api-scenarios.json \\
    DB (H2) & DbTest & db-scenarios.json \\
    DB (SQL Server) & SqlServerDbTest & db-sqlserver-scenarios.json \\
    E2E & E2EOrderFlowTest & e2e-scenarios.json \\
    \bottomrule
  \end{tabular}
\end{table}

\begin{table}[H]
  \centering
  \caption{TestNG suite grupları (ui-regression.xml)}
  \begin{tabular}{@{}ll@{}}
    \toprule
    \textbf{Grup} & \textbf{İçerik} \\
    \midrule
    01-Public-Account-Flows & Login, Register, ForgotPassword \\
    02-My-Account & MyAccount, MyAccountSettings, Newsletter \\
    03-Storefront & WishList, Search, Cart, Checkout \\
    04-API-Layer & ApiTest \\
    05-DB-Layer & DbTest, SqlServerDbTest \\
    06-E2E-Cross-Layer & E2EOrderFlowTest \\
    \bottomrule
  \end{tabular}
\end{table}


\chapter{Ek B: Örnek Konfigürasyon}
\begin{lstlisting}[language=bash,caption={config.example.properties özeti}]
base.url=https://ecommerce-playground.lambdatest.io/...
api.base.url=https://reqres.in
ai.provider=groq
ai.model=llama-3.1-8b-instant
run.ui.layer=true
run.api.layer=true
run.db.layer=true
run.e2e.layer=true
\end{lstlisting}

\clearpage
\section*{KAYNAKLAR}
\addcontentsline{toc}{section}{KAYNAKLAR}

\begin{thebibliography}{99}

\bibitem{istqb}
International Software Testing Qualifications Board (ISTQB).
\textit{Certified Tester Foundation Level (CTFL) v4.0}.
ISTQB, 2023.
Erişim: \url{https://istqb.org/certifications/certified-tester-foundation-level-ctfl-v4-0/}

\bibitem{owasp-wstg}
OWASP Foundation.
\textit{Web Security Testing Guide (WSTG)}.
OWASP, 2024.
Erişim: \url{https://owasp.org/www-project-web-security-testing-guide/}

\bibitem{owasp-auth}
OWASP Foundation.
\textit{Authentication Cheat Sheet}.
OWASP Cheat Sheet Series, 2024.
Erişim: \url{https://cheatsheetseries.owasp.org/cheatsheets/Authentication_Cheat_Sheet.html}

\bibitem{owasp-input}
OWASP Foundation.
\textit{Input Validation Cheat Sheet}.
OWASP Cheat Sheet Series, 2024.
Erişim: \url{https://cheatsheetseries.owasp.org/cheatsheets/Input_Validation_Cheat_Sheet.html}

\bibitem{selenium}
Selenium Project.
\textit{Selenium WebDriver Documentation}.
2024.
Erişim: \url{https://www.selenium.dev/documentation/}

\bibitem{testng}
TestNG Project.
\textit{TestNG Documentation}.
2024.
Erişim: \url{https://testng.org/}

\bibitem{pom}
Fowler, M.
\textit{Page Object}.
martinfowler.com, 2013.
Erişim: \url{https://martinfowler.com/bliki/PageObject.html}

\bibitem{maven}
Apache Software Foundation.
\textit{Apache Maven Documentation}.
2024.
Erişim: \url{https://maven.apache.org/guides/}

\bibitem{webdrivermanager}
Bonigarcia, B.
\textit{WebDriverManager Documentation}.
2024.
Erişim: \url{https://bonigarcia.dev/webdrivermanager/}

\bibitem{restassured}
REST Assured Project.
\textit{REST Assured}.
2024.
Erişim: \url{https://rest-assured.io/}

\bibitem{jackson}
FasterXML.
\textit{Jackson Databind}.
GitHub, 2024.
Erişim: \url{https://github.com/FasterXML/jackson-databind}

\bibitem{extentreports}
Extent Reports.
\textit{ExtentReports 5 --- Java Documentation}.
2024.
Erişim: \url{https://www.extentreports.com/docs/versions/5/java/index.html}

\bibitem{lambdatest}
LambdaTest.
\textit{E-Commerce Playground}.
2024.
Erişim: \url{https://ecommerce-playground.lambdatest.io/}

\bibitem{reqres}
ReqRes.
\textit{ReqRes --- A Hosted REST-API Ready to Respond to Your Requests}.
2024.
Erişim: \url{https://reqres.in/}

\bibitem{chartjs}
Chart.js Contributors.
\textit{Chart.js Documentation}.
2024.
Erişim: \url{https://www.chartjs.org/docs/latest/}

\bibitem{groq}
Groq Inc.
\textit{Groq API Documentation}.
2024.
Erişim: \url{https://console.groq.com/docs}

\bibitem{github-actions}
GitHub.
\textit{GitHub Actions Documentation}.
2024.
Erişim: \url{https://docs.github.com/en/actions}

\end{thebibliography}


\end{document}
